\chapter{Preface}
\label{chap:preface}

The purpose of this thesis is to produce a tutorial. A working system was built in the course of it, and that system is described here in full, though the intended result of the work is not only the system itself but also the documented procedure that leads to it. The tutorial form is therefore an objective rather than a stylistic decision, and it governs how the material is arranged. Every step is given in the order in which it must be performed, so that a reader starting from a bare board and an empty toolchain can reproduce the whole result.

The system to be built is a complete video system on a single \term{SoC FPGA}, with acquisition carried out in hardware, and processing carried out in software, on the processor system integrated in the same device.

\section{The hardware provided}

The assigned platform is a \term{Terasic DE1-SoC} development board, built around an \term{Intel Cyclone V SoC}. This device is an \term{FPGA} that integrates a processor: on the same die sit the reconfigurable logic and a hard processor system (\term{HPS}), a dual-core \term{ARM Cortex-A9} with its own \term{DDR3} memory controller, joined by dedicated \term{HPS-to-FPGA} bridges. Custom logic can therefore be addressed by software as ordinary memory-mapped registers and can in turn reach the processor's own \term{DDR3} memory directly. Almost every design decision described in the following chapters follows from this arrangement.

The image sensor is an \term{OV7670} camera module, a 640$\times$480 \term{CMOS} rolling-shutter device capable of 30 frames per second and able to code each pixel as Bayer mosaic, \term{RGB565}, \term{RGB555}, \term{RGB444} or \term{YUV 4:2:2}. Its transmission protocol is strictly unidirectional, with no handshaking and no flow control, so pixels leave the sensor on an 8-bit parallel bus timed by a pixel clock and delimited by line and frame signals. The receiver then has to follow the sensor's timing rather than ask it for data. The pixel clock is not free-running, since it is derived from a master clock\footnote{TODO: footnote text was lost when the source was copy-pasted -- please supply it.} which must be supplied from outside. The sensor also exposes an \term{I\textsuperscript{2}C}-like control bus, which is used to set up resolution, output format and the other imaging parameters at run time by writing its configuration registers. Reading the video stream from a such device is therefore not a trivial connection to a well-known interface, but the construction of a custom one, whose logic interprets its timing and forwards the data to the \term{HPS} through the \term{FPGA} fabric.

During development it became clear that the supplied jumper leads, about ten centimetres long, degraded the video signal, and that the twenty-centimetre ones, needed to move the camera freely, corrupted it altogether. For this reason, a pair of small \term{PCB}, one at each end of the cable, was designed and built to drive the data line. Generating the camera master clock\footnote{TODO: footnote text was lost when the source was copy-pasted -- please supply it.} in the \term{FPGA} was weighed at that point and set aside, since it would have called for an additional line buffer, and a cheaper clock generator sitting next to the camera served instead.

\section{Requirements and constraints}

Two requirements shaped the development. The first, set by the project, is to relieve the processor of the image data transfer, or at least to involve it as little as possible. Its computing resources are meant for computing, not for carrying frames from the acquisition logic into memory: keeping a processor busy with a sustained transfer that involves no computation at all wastes both \term{CPU}'s cycles and energy.

The second is self-imposed and softer, that no step of the procedure depends on a paid license. Every tool used here, from the \term{FPGA} design software to the compiler and the host utilities, is free of charge or open source, so that anyone taking this work further can go from a bare board to a running system without buying anything.

\section{Outline of the implementation}

\subsection{The hardware}

Given the first requirement, the natural answer is a \term{DMA} unit, able to carry the data on its own, straight from the acquisition logic into the processor's \term{DDR3} memory. All the software has to supply is the address at which the incoming frame is to be stored.

The board provided, and the \term{Cyclone V SoC} at its center, offers no such unit ready in the \term{HPS}. The \term{DMA} must therefore be implemented in the \term{FPGA} fabric and attached to the \term{HPS}'s \term{DDR3} controller through the \term{HPS-to-FPGA} bridges it exposes.

This can be accomplished with the manufacturer's own design software, \term{Quartus Prime}, which includes \term{Platform Designer}, a tool for composing a system out of hardware blocks.\footnote{TODO: footnote text was lost when the source was copy-pasted -- please supply it.} They may be drawn from the vendor's own library or written for the occasion: each one must only present a compatible interface. \term{Platform Designer} then generates the interconnect that joins them together, in the form of hardware description language. Intel supplies a \term{Modular Scatter-Gather DMA}, or \term{mSGDMA}, which fits the requirement exactly and is the unit used throughout this work. The remaining modules that complete the capture path are realized as \term{Platform Designer} components too, custom ones when the library held nothing that did the job.

The generic part of the hardware, essentially everything outside the video path, follows two reference designs, a published design for the \term{DE1-SoC} [UPM-DE1] and Intel's \term{Golden Hardware Reference Design} [GHRD], the latter targeting a different board built around a similar \term{Cyclone V SoC}. Neither could be reused as it stood, one because it targets another board and the other because it has no video path, but together they cover the obligatory groundwork, how the \todo{word missing in the original source -- fill in on reread} is instantiated and clocked, how reset is distributed and how the pins of the board are assigned. That part of the design admits one correct solution, so there was nothing to invent there.

\subsection{The software}

The software side\footnote{TODO: footnote text was lost when the source was copy-pasted -- please supply it.} took the form of two projects, kept in separate repositories,\footnote{TODO: footnote text was lost when the source was copy-pasted -- please supply it.} \sysname{FPGAlix}, which runs under embedded \term{Linux} on the \term{HPS}, and \sysname{FPGAsteel}, which runs bare metal on the same processor. The two environments differ in what they make affordable. For bringing the board up, \term{Linux} offers a safe road, since the boot chain, the drivers and the capture framework are code that has been made to work before, so when the system does not come up the fault lies in how it has been configured rather than in the software itself. For the application it supplies device drivers, memory management, cache management, a standard video capture framework and tested libraries, which shorten the work, and it allows the program to be debugged remotely over \term{Ethernet} with |gdb|, while bare metal has to be followed through \term{USB JTAG}. What it costs is distance from the hardware and timing that is not entirely predictable. Bare metal removes the abstraction and the overhead and gives complete control of the entire system, but everything must be written from nothing, the processing code among it, and the problems the operating system hides come back into view.

From the hardware point of view the two projects differ slightly. In both a single frame can be read out through the debugger and examined on the development host, while the display of a continuous video stream requires a dedicated output path, and the two projects adopt different ones. \sysname{FPGAlix} transmits the frames over \term{Ethernet}, relying on the interface and the network stack provided by the operating system, at the cost of a significant processor load, due to encoding and to the splitting of the stream into network packets. \sysname{FPGAsteel} drives the \term{VGA} port available on the board,\footnote{TODO: footnote text was lost when the source was copy-pasted -- please supply it.} since implementing a network driver and a \term{TCP/IP} stack in bare metal would have required an effort disproportionate to the aim of the project.

\section{Organization of the document}

The two projects just introduced, \sysname{FPGAlix} and \sysname{FPGAsteel}, are described in the two parts that follow in the tutorial form the aim of the work requires, with every step given in the order in which it must be carried out. A working familiarity with \term{Quartus Prime}, its \term{Platform Designer}, with \term{C++} and with the \term{Linux} command line is assumed. A reader who lacks the first two will find in [UPM-DE1] a guided introduction which, although written for an older release of the tools, still describes the \term{Quartus} side closely enough to be followed. What the present document explains is the principle of operation rather than the code. The sources themselves can be cloned from the repositories: the hardware is described in \term{VHDL}, the software is written in \term{C++} and built with |make|.

\textbf{Part I} covers \sysname{FPGAlix}. It defines the first version of the hardware, brings up the operating system on the HPS and closes with one example of video processing. A single example is enough here, because under the operating system the software is divided into modules with well-defined boundaries, each of them behaving in a familiar and predictable way, so further examples would add length without adding understanding.

\textbf{Part II} covers \sysname{FPGAsteel}. On the hardware side it adds the \term{DMA} that feeds the \term{VGA} video output, and on the software side it offers several examples. Bare-metal code is tailored to a single system and very little of it is standard since it heavily depends on the choices of whoever writes it. The examples therefore proceed by steps, each one adding a single capability to the one before it:

\begin{enumerate}
    \item prints a message on the serial console,
    \item establish an \term{I\textsuperscript{2}C} communication for configuring the camera,
    \item forwards the incoming video stream from the camera to \term{VGA} output,
    \item the same path then is tested with the caches enabled,
    \item brings up the second core,
    \item implements an image-processing chain, that runs on both cores, for locating the boundaries and the outline of a bright object against a dark background.
\end{enumerate}

\reread{Part III gathers the appendices: the schematics and the bill of materials of the camera adapter board, and the bibliography, which lists the project repositories, the external documents referred to throughout and the artefacts produced by the work. The three parts together are meant to be sufficient to rebuild the whole system starting from a bare board, and each of the first two can also be read on its own by a reader interested only in the \term{Linux} side or only in the bare-metal one.}
